%%%%%%%%%%%%%%%%%%%%%%%%%%%%%%%%%%%%%%%%%%%%%%%%%%%%%%%%%
%                                                       %
%  DOCUMENTO Nº5: ANEXOS AL MANUAL DEL ADMINISTRADOR    %
%                                                       %
%%%%%%%%%%%%%%%%%%%%%%%%%%%%%%%%%%%%%%%%%%%%%%%%%%%%%%%%


\documentclass[a4paper,12pt]{article}
%Fichero con configuración opcional:
%
% Paquetes que pueden serte de utilidad (rec = recomendado, opc = opcional)
%
\usepackage{fancyhdr}          % (rec)  permite cambiar varios par�metros de las cabeceras y pi�s de p�gina
% \usepackage{fancyvrb}		% (rec) permite cambiar parámetros del texto
% literal (entorno verbatim)
\usepackage{courier}           % (opc)  usa esta fuente por defecto
%\usepackage{serif}
\usepackage[spanish]{babel}   % (rec)  da soporte para castellano a LaTeX
\usepackage[utf8]{inputenc}    % (rec) soporte para tildes
%\usepackage{setspace}          % (opc)  permite cambiar el espaciado entre l�neas
%\usepackage{longtable}         % (opc)  permite que las tablas ocupen varias p�ginas
%\usepackage{lscape}            % (opc)  permite el uso del comando \landscape, para poner algo apaisado
%\usepackage{color}             % (opc)  varios comandos relativos al color (como \color)
\usepackage{colortbl}           % (opc) permite usar color en las tablas
% \usepackage{rotating}          % (opc)  permite rotar PSs y EPSs
% \usepackage{textcomp}          % (opc)  permite incluir el s�mbolo del euro,
% con \texteuro
\usepackage{gensymb} 		% para meter el s\'imbolo del ohmio
%\usepackage{mathabx}		% para meter el simbolo mayor menor
%\usepackage{minitoc}           % (opc)  permite incluir ToCs (�ndice de materias) para cada cap�tulo
%\usepackage{epsf}              % (opc)  permite ciertas manipulaciones a EPSs
\usepackage[absolute]{textpos} % (rec)  permite posicionado arbitrario de texto (necesario para la portada)
%\usepackage{srcltx}            % (opc)  permite pasar del .dvi al .tex
% \usepackage{marvosym}
\usepackage[breaklinks=true]{hyperref} % (opc) hace que las referencias cruzadas se conviertan en enlaces
\hypersetup{
    bookmarks=true,         % show bookmarks bar?
    unicode=false,          % non-Latin characters in Acrobat’s bookmarks
    pdftoolbar=true,        % show Acrobat’s toolbar?
    pdfmenubar=true,        % show Acrobat’s menu?
    pdffitwindow=false,     % window fit to page when opened
    pdfstartview={FitH},    % fits the width of the page to the window
    pdftitle={MASER},    % title
    pdfauthor={Digna Gonz\'alez, Ivan Iglesisas},     % author
    pdfsubject={MASER},   % subject of the document
%     pdfcreator={Creator},   % creator of the document
%     pdfproducer={Producer}, % producer of the document
%     pdfkeywords={keyword1} {key2} {key3}, % list of keywords
    pdfnewwindow=true,      % links in new window
    colorlinks=true,       % false: boxed links; true: colored links
    linkcolor=red,          % color of internal links
    citecolor=green,        % color of links to bibliography
    filecolor=magenta,      % color of file links
    urlcolor=cyan           % color of external links
}


\usepackage{titleref}           % (opc) permite hacer referencia a otra sección
% mostrando su título
\usepackage{graphicx}
% \usepackage{multirow}
% \usepackage{setspace}
\usepackage{url}     % Enlaces web
% \usepackage{tabularx} %Ajuste automático del ancho de columnas
% \usepackage{longtable} %Tablas de más de una página, por ejemplo para los
% acrónimos
% \usepackage{listings} %Mostrar código en el texto, por ejemplo comandos de
% consola
% \usepackage{cmap}   %Hacer que el pdf sea searchable

% \definecolor{gray97}{gray}{.97}
% \definecolor{gray75}{gray}{.75}
% \definecolor{gray45}{gray}{.45}
% \definecolor{BitterSweet}{cmyk}{0,0.83,1,0.70}
% \definecolor{BlueViolet}{cmyk}{0.86,0.91,0,0.04}
%  
% \lstset{ frame=Ltb,
%      framerule=0pt,
%      aboveskip=0.5cm,
%      framextopmargin=3pt,
%      framexbottommargin=3pt,
%      framexleftmargin=0.4cm,
%      framesep=0pt,
%      rulesep=.4pt,
%      backgroundcolor=\color{gray97},
%      rulesepcolor=\color{black},
% 	inputencoding=utf8,
% 	extendedchars=\true,
%      %
%      stringstyle=\ttfamily,
%      showstringspaces = false,
%      basicstyle=\small\ttfamily,
%      commentstyle=\color{gray45},
%      keywordstyle=\bfseries,
%      %
%      numbers=left,
%      numbersep=15pt,
%      numberstyle=\tiny,
%      numberfirstline = false,
%      breaklines=true,
%    }
%  
% % minimizar fragmentado de listados
% \lstnewenvironment{listing}[1][]
%    {\lstset{#1}\pagebreak[0]}{\pagebreak[0]}
%  
% \lstdefinestyle{consola}
%    {basicstyle=\scriptsize\bf\ttfamily,
%     backgroundcolor=\color{gray75},
%    }
% 
% \lstdefinestyle{archivo}
%    {basicstyle=\scriptsize\ttfamily,
%     backgroundcolor=\color{gray97},
%    }

%\usepackage{urlbst} %Soporte para url en bibtex
%
% Settings para los m�rgenes. Descomenta y modifica si sabes lo que haces. N�tese
% que a los valores dados se les a�ade una pulgada extra. Los valores dados son los
% predeterminados para papel A4 y el estilo itsas_pfc.cls.
%
%\setlength{\oddsidemargin}{10pt}     % m�rgen izquierdo para p�ginas impares (izquierda)
%\setlength{\evensidemargin}{52pt}    % m�rgen izquierdo para p�ginas pares (derecha)
%\setlength{\textwidth}{390pt}        % anchura del cuerpo de texto

%
% Recomendado para mejorar la colocaci�n autom�tica de las figuras.
% (tomado de http://dcwww.camp.dtu.dk/~schiotz/comp/LatexTips/LatexTips.html#captfont)
%
\renewcommand{\topfraction}{0.85}
\renewcommand{\textfraction}{0.1}
\renewcommand{\floatpagefraction}{0.75}

%
% Espacio entre el borde superior de la p�gina y donde comienza el texto (ah� van las
% cabeceras). LaTeX se queja si usamos el paquete fanchyhdr y headheight es menor de 15pt
%
\headheight 16pt

%
% Para el paquete textpos (usado para la portada)
%
\setlength{\TPHorizModule}{\paperwidth}
\setlength{\TPVertModule}{\paperheight}
\newcommand{\tb}[4]{\begin{textblock}{#1}[0.5,0.5](#2,#3)\begin{center}#4\end{center}\end{textblock}}
% 
%
% \newcommand{\respuesta}[1]{\setlength{\parindent}{0pt}\colorbox[gray]{0.87}{
% \scriptsize\texttt{\BUseVerbatim{#1}}}\setlength{\parindent}{1cm}}
% 
% \newcommand{\enunciado}[1]{\textit{\textbf{#1}}}



%
% Aqu� puedes definir tus comandos.
% 
% \newcommand{cmd}[args]{def}
%
% cmd  = el comando a definir (p.e. \cadena)
% args = el n�mero de argumentos
% def  = la definici�n, sustituyendo #1, #2... por el primer, segundo... argumento
%
% Por ejemplo:
%
% \newcommand{\agua}[1]{H\ensuremath{_#1}O}
%
% Cada vez que escribamos "\agua{33}", en el output saldr�: "H33O" (con el 33 como sub�ndice)
%
% 
% %\newcommand{\algo}{algo}
% 
% \newcommand{\todolist}[1]{
%  \marginpar{
%   \fbox{
%    \begin{minipage}{3cm}
%     \tiny{
%      \begin{list}{$\bullet$}
%       {
%       \textsc{\textbf{To do: }}\vspace*{-0.2cm}
%       \setlength\labelwidth{0.1cm}
%       \setlength\itemindent{0cm}
%       \setlength\leftmargin{0.05cm}
%       \setlength\parsep{0cm}
%       \setlength\itemsep{0.05cm}
%       }
%       #1
%      \end{list}
%     }
%    \end{minipage}
%   }
%  }
% }

%
% Aqu� puedes instruir a LaTeX de por d�nde cortar las palabras que �l autom�ticamente
% no sepa. P.e., para cortar "gnomonly" solo por donde se se�ala con guiones (-).
%
% \hyphenation{gno-mon-ly} 
 
%
% Que las primeras p�ginas sean numeradas con n�meros romanos.
% M�s adelante se cambiar� de nuevo a ar�bicos.
%
% \pagenumbering{Roman}


\begin{document}
%Da valores a ciertas variables:
%
% Este fichero contiene una lista de nombres (variables) internos
% de LaTeX, a los que puedes cambiar el nombre. Por ejemplo, puedes
% hacer que los cap�tulos se llamen "Secci�n" en vez de "Cap�tulo"
%
\renewcommand\bibname{Bibliograf\'{\i}a}        % as� el nombre de la secci�n Bibliograf�a ser� "Bibliograf�a".
\newcommand{\myname}{Digna González Otero}                 % nombre del autor.
%\newcommand{\myboss}{José Daniel Gutiérrez Porset}                 % nombre del supervisor.
\newcommand{\thesistitle}{\uppercase{Construcción e implantación de herramientas para la gestión colaborativa del conocimiento en el GSC}}        % t�tulo del trabajo.
\newcommand{\worktype}{Proyecto Fin de Carrera} % tipo de trabajo.
\newcommand{\upv}{Utils/ehu_logo.eps}          % fichero con el logo (p.e. para la portada).
\newcommand{\ibi}{Utils/logo_ingenieria.eps}          % fichero con el logo (p.e. para la portada).
\newcommand{\logoupv}{Utils/ehu_logo.png}          % fichero con el logo (p.e. para la portada).
\newcommand{\logoibi}{Utils/logo_ingenieria.png}          % fichero con el logo (p.e. para la portada).
%\renewcommand{\figurename}{xxx}                % nombre a pie de figura (xxx 1: bla-bla-bla).
%\renewcommand{\listfigurename}{yyy}            % nombre del �ndice de figuras.
\renewcommand{\tablename}{Tabla}
\renewcommand{\listtablename}{Índice de tablas}

    \makeatletter
    \def\thebibliography#1{\section*{REFERENCIAS\@mkboth
      {REFERENCIAS}{REFERENCIAS}}\list
      {[\arabic{enumi}]}{\settowidth\labelwidth{[#1]}\leftmargin\labelwidth
	\advance\leftmargin\labelsep
	\usecounter{enumi}}
	\def\newblock{\hskip .11em plus .33em minus .07em}
	\sloppy\clubpenalty4000\widowpenalty4000
	\sfcode`\.=1000\relax}
    \makeatother

%Página de portada:
%
% Este fichero produce la primera p�gina del PFC/Tesis, con el
% t�tulo, tu nombre, el nombre del supervisor, etc.
%
% La mayor�a, si no todo el contenido de esta p�gina est� incluido a trav�s 
% de comandos (p.e. \thesistitle) que han sido definidos en Config/pfc_options.tex
%
% Editar a voluntad.
%

\thispagestyle{empty} % no imprimir ni n�mero, ni cabecera ni pi� de p�gina

%
% Usa \tb para colocar varios items en la p�gina. Uso:
%
% \tb{w}{h}{v}{t}
%
% donde:
%
% w = ancho de la caja con el texto (1.0 = ancho de p�gina)
% h = posici�n horizontal del centro de la caja de texto (0.0 = izquierda, 1.0 = derecha)
% v = posici�n vertical del centro de la caja de texto (0.0 = arriba, 1.0 = abajo)
% t = texto a incluir en la caja de texto
%
\tb{0.4}{0.25}{0.1}{\includegraphics[width=0.8\textwidth]{\logoupv}} % logo UPV/EH
\tb{0.6}{0.7}{0.1}{\includegraphics[width=0.65\textwidth]{\logoibi}} % logo UPV/EHU
%\includegraphics[width=0.5\textwidth]{Utils/ehu_logo.png}
\tb{0.8}{0.5}{0.17}{\rule{\columnwidth}{0.025cm}}
\tb{0.8}{0.50}{0.35}{\huge \thesistitle}                 % t�tulo
%%\tb{0.8}{0.50}{0.40}{\Large \worktype}                        % si es PFC o Tesis
\tb{0.5}{0.50}{0.55}{\huge Documento nº 7:}
\tb{0.8}{0.50}{0.60}{\large \textbf{ANEXOS AL MANUAL DEL ADMINISTRADOR}}
\tb{0.8}{0.50}{0.75}{ALUMNA:	Digna María González Otero}
\tb{0.8}{0.50}{0.78}{DIRECTOR:	José Daniel Gutiérrez Porset}
\tb{0.8}{0.80}{0.9}{Bilbao, Marzo de 2009}


\ \clearpage                       % terminar la p�gina aqu�%
\thispagestyle{empty} \ \clearpage % p�gina en blanco

% CONFIG: Inicializar en 1 el n mero de p gina
\input{Config/begin.tex}
% CONFIG:  ndice de materias (ToC)
%\dominitoc        % que cada cap�tulo tenga su ToC (necesita el paquete "minitoc" antemencionado)
\tableofcontents  % insertar ToC en este punto
\newpage
\listoffigures    % insertar lista de Figuras (LoF) en este punto (opcional)
% \newpage
% \listoftables     % insertar lista de Tablas (LoT) en este punto (opcional)

\cleardoublepage

% CONFIG: Definir el estilo de las cabeceras/pi s de p gina.
\pagestyle{fancy}                                % elegir este estilo de cpps (recomendado)
\fancyhf{}                                       % borra el estilo anterior para cpps, para luego redefinirlos
\fancyfoot[C]{\textbf{\thepage}}             % Cabecera: n�mero de p�gina en negrita.
\fancyhead[R]{\nouppercase{\leftmark}}          % Cabecera: incluye informaci�n del nivel superior (Cap�tulo)
\renewcommand{\headrulewidth}{0.5pt}             % Cabecera: subraya la cabecera (fijar en "0pt" si no se desea).
\renewcommand{\footrulewidth}{0.5pt}               % Pi�: subraya el pie de p�gina (fijar en "0pt" si no se desea).






% CONFIG: Estilo de los cap tulos
\setcounter{page}{1}    % empezar a contar de nuevo desde 1 las p�ginas.
\pagenumbering{arabic}  % utilizar n�meros �rabes de nuevo.
\setlength{\parskip}{0.5cm} %Separación entre párrafos

% Usa \tocchapter en vez de \chapter, para usar cap�tulos
% bien formateados:
%\newcommand{\tocchapter}[1]{\cleardoublepage\chapter{#1}\minitoc\newpage}

% CONTENIDO: Tras esto puedes incluir todos los cap tulos/secciones que desees.
\section{Introducción}
\label{introduccion}

Un m\'aser es un \textbf{dispositivo que produce ondas electromagn\'eticas coherentes} mediante la amplificaci\'on por la emisi\'on estimulada de radiaci\'on \cite{maserDefinition}. El t\'ermino se defini\'o como acr\'onimo de \textit{Microwave Amplification by Stimulated Emission of Radiation}.

Como su nombre indica, su funcionamiento está basado en el fenómeno de emisión estimulada de radiación, enunciado por Albert Einstein en 1916.

Una de las principales aplicaciones de los m\'aseres es como relojes at\'omicos de gran precisi\'on, pero tambi\'en se usan como amplificadores, principalmente en radiotelescopios. En secciones posteriores de este documento profundizaremos en las aplicaciones de este tipo de dispositivos.

\subsection{Historia}

El funcionamiento del m\'aser est\'a basado en el fen\'omeno de emisi\'on estimulada de radiaci\'on enunciado por Albert Einstein en 1916. 

En mayo de 1952 Nikolay Basov y Alexander Prokhrov del Instituto de F\'isica de Lebedev describieron el principio del m\'aser en la Academia de Ciencias de la URSS. M\'as tarde, en octubre de 1954  publicaron los resultados de su investigaci\'on.

De forma independiente Charles H. Townes, J. P. Gordon y H. J. Zeiger fabricaron en 1953 en la Universidad de Columbia el primer m\'aser. Este dispositivo utilizaba un haz de mol\'eculas de amon\'iaco para producir la amplificaci\'on de microondas a la frecuencia de 24 GigaHerzios.

M\'as tarde Townes trabaj\'o con Arthur L. Schawlow para describir el principio del m\'aser \'optico o l\'aser, y en 1958 hicieron varias publicaciones sobre este tema, aunque por aqu\'el entonces no siguieron con las investigaciones en ese campo.

En 1959 Townes y Schawlow obtuvieron una patente del m\'aser. Este dispositivo se usaba para amplificar se\~nales de radio y como detector ultrasensible para investigaci\'on espacial.

En 1964 Townes, Basov y Prokhorov recibieron el premio Nobel de F\'isica por su trabajo en este campo. A Townes se le adjudic\'o la mitad del premio y la otra mitad se reparti\'o entre Basov y Prokhorov.

\begin{figure}[h]
\begin{center}$
\begin{array}{ccc}
\includegraphics[width=0.25\textwidth]{./Utils/townes.jpg} &
\includegraphics[width=0.25\textwidth]{./Utils/basov.jpg} &
\includegraphics[width=0.25\textwidth]{./Utils/prokhorov.jpg}
\end{array}$
\end{center}
\caption{Townes, Basov y Prokhorov, de izquierda a derecha}
\end{figure}

En 1965 se descrubri\'o el primer m\'aser astron\'omico en la nebulosa de Ori\'on. Era un m\'aser de hidroxilo, que emite radiaci\'on en varias longitudes de onda, aunque las m\'as utilizadas son las de 18cm (1,7GHz de frecuencia).
\subsection{Terminolog\'ia}

Aunque el t\'ermino \textit{MASER} se difini\'o como acr\'onimo de \textbf{Microwave Amplification by Stimulated Emission of Radiation}, con el tiempo y el avance de la tecnolog\'ia este t\'ermino se ha ido quedando obsoleto, y se ha empezado a llamar al dispositivo \textit{m\'aser} en min\'usculas, y ya no en forma de acr\'onimo. 

El motivo por el que la definici\'on ya no es precisa es debido a que el rango de frecuencias de funcionamiento de los m\'aseres actuales no est\'a limitado al de microondas, sino que se extiende en el espectro. Por este motivo el f\'isico Charles H. Townes, quien le hab\'ia dado el nombre original junto con sus colaboradores, sugiri\'o sustituir la palabra ``microondas'' por ``molecular'' en la definici\'on. Otros cient\'ificos, sin embargo, alegaron que con este intento de extender el acr\'onimo Charles H. Townes estaba intentando darle m\'as importancia a su descubrimiento y aumentar su reputaci\'on en el mundo cient\'ifico.

Cuando se desarroll\'o el l\'aser en 1957, Townes, Schawlow y sus colegas de Bell Labs impulsaron la adopci\'on del t\'ermino \textit{m\'aser \'optico}, pero \'este fue sustituido por \textit{l\'aser} (Light Amplification by Stimulated Emission of Radiation), nombre defendido por Gordon Gould. De hecho Gould propuso un nombre diferente para los dispositivos que emit\'ian en cada regi\'on del espectro. As\'i, habr\'ia \textit{grasers} (``gamma ray lasers''), \textit{xasers} (``x-ray lasers''), \textit{uvasers} (``ultraviolet lasers''), \textit{lasers} (``visible lasers''), \textit{iraser} (``infrared lasers''), \textit{masers} (``microwave masers'') y \textit{rasers} (``RF masers''). La mayor parte de estos nombres, sin embargo, nunca llegaron a adoptarse, y todos se han quedado obsoletos excepto \textit{\textbf{m\'aser}} y \textit{\textbf{l\'aser}}.

Actuamente se denomina \textit{l\'aser} al dispositivo que emite en la regi\'on del espectro delimitada por los rayos-X y los infrarrojos, y \textit{\textbf{m\'aser}} al que \textbf{emite en la regi\'on de microondas e inferiores}. 

\clearpage
\newpage
\section{Instalación y configuración de Kubuntu}
\label{instKubuntu}

En esta sección se describe el proceso de instalación y configuración de Kubuntu 8.04 (Hardy Heron).

\subsection{Instalación de Kubuntu}

En primer lugar hay que configurar el ordenador para que arranque desde la unidad de CD. A continuación se introduce el disco de Kubuntu en el lector y se reinicia el equipo para empezar el proceso de instalación.

Los pasos a realizar para llevar a cabo la instalación son los siguientes:

\begin{enumerate}
 \item Selección de idioma para el proceso de instalación (español, en este caso)
\begin{figure}[htb!]
 \centering
 \includegraphics[width=0.8\textwidth]{Utils/instalacion1.png}
 % instalacion1.png: 1179668x1179666 pixel, 0dpi, infxinf cm, bb=
 \caption{Selección de idioma de la instalación}
 \label{fig:instalacion1}
\end{figure}

\item Seleccionar la opción ``Instalación de Kubuntu'' en el menú

\item Elección de ubicación (España en el caso que nos ocupa)

\begin{figure}[htb!]
 \centering
 \includegraphics[width=0.8\textwidth]{Utils/instalacion4.png}
 % instalacion1.png: 1179668x1179666 pixel, 0dpi, infxinf cm, bb=
 \caption{Selección de la ubicación}
 \label{fig:instalacion4}
\end{figure}

\item Selección del modelo de teclado (español).

\item Detección de hardware, descarga de componentes adicionales necesarios y configuración de red (procesos automáticos).

\item Particionado del sistema. Es posible mantener particiones que contengan otros sistemas operativos, aunque en este caso se destinará toda la capacidad del disco a Kubuntu. Para ello se elegirá el particionado manual y se eliminarán todas las particiones existentes. A continuación se crearán tres particiones utilizando el gestor de particiones mostrado en la Figura \ref{fig:particionado}:
\begin{itemize}
 \item Partición de 19GB con sistema de ficheros ext3 y punto de montaje en ``/''.
 \item Partición de 60Gb con sistema de ficheros ext3 y punto de montaje en ``/home''.
 \item Partición de  1GB de tipo swap (intercambio).
\end{itemize}

\begin{figure}[htb!]
 \centering
 \includegraphics[width=0.8\textwidth]{Utils/particionado.png}
 % instalacion1.png: 1179668x1179666 pixel, 0dpi, infxinf cm, bb=
 \caption{Particionado del sistema}
 \label{fig:particionado}
\end{figure}

\item Creación del sistema de archivos ext3 e instalación del sistema base Kubuntu. Descarga del resto de paquetes necesarios a disco, incluyendo el soporte para el idioma español. (Proceso automático.)

\item Introducción del nombre del administrador del sistema, nombre de usuario y contraseña y nombre del equipo.

\item Fin de la instalación inicial de Kubuntu. Es necesario extraer el disco del lector de CD y reiniciar el sistema. Este primer arranque será bastante más lento de lo habitual ya que se descargarán e instalarán numerosos paquetes.
\end{enumerate}

\newpage
Una vez hecho introducidos todos los datos y guardados los cambios la conexión a
Internet ya estará activa.
\clearpage
\newpage
\section{Instalación de DokuWiki}
\label{instDokuwiki}

Esta sección describe la instalación de la plataforma para la creación de wikis Dokuwiki y de los elementos necesarios para su funcionamiento.

\begin{enumerate}
\item Instalación del servidor web Apache y los módulos php necesarios:

\begin{listing}[style=consola, numbers=none]
# aptitude install apache2 php5 php5-apache2-mod-bt
\end{listing}

\item Si se va a llevar a cabo una instalación limpia (sin partir de ninguna plantilla), descarga de la última versión de Dokuwiki de la página oficial del proyecto \cite{dokuwiki}.

\item Descompresión del paquete de DokuWiki (ya sea el descargado de la página o el la plantilla) y copia del mismo en el servidor (directorio \texttt{/var/www/dokuwiki}).

\item Modificar la nueva carpeta para que pertenezca al usuario \textit{www-data} (usuario utilizado por el servidor web):
\begin{listing}[style=consola, numbers=none]
# chown -R www-data /var/www/dokuwiki
\end{listing}

\item Habilitar el módulo \textit{mod\_rewrite} de Apache:
\begin{listing}[style=consola, numbers=none]
# a2enmod rewrite
# /etc/init.d/apache2 restart
\end{listing}

\item Asignar al parámetro \textit{AllowOverride} el valor \textit{all} en el archivo \textit{default} de Apache:
\begin{listing}[style=consola, numbers=none]
# nano /etc/apache2/sites-available/default
\end{listing}

\item Establecer \texttt{/var/www/dokuwiki} como directorio base de Apache modificando el parámetro \textit{DocumentRoot} del archivo 000-default.
\begin{listing}[style=consola, numbers=none]
# nano /etc/apache2/sites-enabled/000-default
\end{listing}

\item Configurar la página de inicio del servidor web estableciendo ``doku.php'' como primera opción en el parámetro DirectoryIndex del archivo \texttt{dir.conf}

\begin{listing}[style=consola, numbers=none]
# nano /etc/apache2/mods-enabled/dir.conf
\end{listing}

\item Acceder a \url{http://localhost/install.php} y rellenar los datos solicitados por el instalador.

\begin{figure}[htb!]
 \centering
 \includegraphics[width=0.8\textwidth]{Utils/instalador_wiki.png}
 % instalacion1.png: 1179668x1179666 pixel, 0dpi, infxinf cm, bb=
 \caption{Instalador web de DokuWiki}
 \label{fig:instalador_wiki}
\end{figure}

\item Una vez finalizada la instalación de DokuWiki, eliminar el archivo \texttt{install.php}:
\begin{listing}[style=consola, numbers=none]
# rm /var/www/dokuwiki/install.php
\end{listing}

\item El wiki ya es accesible de forma local. Para que sea accesible también por otros equipos de la red local habrá que modificar el archivo 000-default para establecer la opción \textit{allow from} a \textit{all}.
\begin{listing}[style=consola, numbers=none]
# nano /etc/apache2/sites-enabled/000-default
\end{listing}
\end{enumerate}


\clearpage
\newpage
\section{Instalación y configuración de Samba}
\label{Samba}

\subsection{Instalación}

Samba es una implementación para GNU/Linux del protocolo SMB (Server Message Block) que permite compartir recursos sobre una red TCP/IP e interconectar ordenadores Windows y Unix.

Para instalar el servidor Samba, basta con instalar el paquete \emph{samba} usando un gestor de paquetes, como apt-get o aptitude.

\begin{listing}[style=consola, numbers=none]
 # aptitude install samba
\end{listing}

El resultado obtenido será el siguiente:

\begin{SaveVerbatim}{codigo}
Leyendo lista de paquetes... Hecho
Creando árbol de dependencias
Leyendo la información de estado... Hecho
Leyendo la información de estado extendido
Inicializando el estado de los paquetes... Hecho
Construir la base de datos de etiquetas... Hecho
Se instalarán los siguiente paquetes NUEVOS:
  samba
0 paquetes actualizados, 1 nuevos instalados, 0 para eliminar y 0 sin actualizar

Necesito descargar 0B/3839kB de ficheros. Después de desempaquetar se usarán 9429kB.
Escribiendo información de estado extendido... Hecho
Preconfigurando paquetes ...
Seleccionando el paquete samba previamente no seleccionado.
(Leyendo la base de datos ...
160960 ficheros y directorios instalados actualmente.)
Desempaquetando samba (de .../samba_3.0.28a-1ubuntu4.4_i386.deb) ...
Configurando samba (3.0.28a-1ubuntu4.4) ...
 * Starting Samba daemons          [ OK ]

Leyendo lista de paquetes... Hecho
Creando árbol de dependencias
Leyendo la información de estado... Hecho
Leyendo la información de estado extendido
Inicializando el estado de los paquetes... Hecho
Escribiendo información de estado extendido... Hecho
Construir la base de datos de etiquetas... Hecho
\end{SaveVerbatim}
\respuesta{codigo}


Si se prefiere, también se pueden descargar los archivos fuente o los binarios de la distribución de la página oficial de Samba \cite{sambadownload} y posteriormente compilarlos e instalarlos siguiendo la documentación sobre la instalación \cite{sambadocu}. Existen manuales en internet que ayudan a llevar a cabo este proceso \cite{sambawiki}.

\subsubsection{Otros paquetes a instalar}

Es interesante instalar también el \textbf{demonio xinetd} (eXtended InterNET Daemon), que es una extensión más segura del servicio de Internet \emph{inetd} que viene instalado por defecto en muchas distribuciones Linux.

Estos demonios gestionan a su vez las conexiones de varios demonios, de forma que se reduce la carga del sistema al ejecutar un único demonio en lugar de todos los que gestiona.

Xinetd contiene mecanismos de control de acceso, permite habilitar los servicios de red en función del tiempo, puede limitar la cantidad de servicios que se ejecutan y contiene un sistema de protección contra escaneos de puertos, entre otras características.

\begin{listing}[style=consola, numbers=none]
# aptitude install xinetd
\end{listing}

Para instalar xinetd habrá que aceptar la desinstalación del paquete openbsd-inetd, correspondiente al demonio inetd.

\begin{SaveVerbatim}{codigo}
Se ELIMINARÁN automáticamente los siguientes paquetes:
  openbsd-inetd
Se instalarán los siguiente paquetes NUEVOS:
  xinetd
Se ELIMINARÁN los siguientes paquetes:
  openbsd-inetd
0 paquetes actualizados, 1 nuevos instalados, 1 para eliminar y 0 sin actualizar
Necesito descargar 0B/137kB de ficheros. Después de desempaquetar se usarán 242kB.
¿Quiere continuar? [Y/n/?] Y
\end{SaveVerbatim}
\respuesta{codigo}

Otro paquete que conviene instalar es \emph{swat}. Este paquete contiene la \textbf{herramienta SWAT} (Samba Web Administration Tool), que permite modificar varios parámetros de configuración del servidor Samba de forma muy sencilla a través de un navegador web.

\begin{listing}[style=consola, numbers=none]
 # aptitude install swat
\end{listing}

Para ver las opciones del comando swat (y comprobar que el paquete se ha instalado correctamente) podemos utilizar el comando \texttt{-?}.
\begin{listing}[style=consola, numbers=none]
 $ swat -?
\end{listing}

\begin{SaveVerbatim}{codigo}
 Usage: swat [OPTION...]
  -a, --disable-authentication       Disable authentication (demo mode)
  -P, --password-menu-only           Show only change password menu

Help options:
  -?, --help                         Show this help message
  --usage                            Display brief usage message

Common samba options:
  -d, --debuglevel=DEBUGLEVEL        Set debug level
  -s, --configfile=CONFIGFILE        Use alternate configuration file
  -l, --log-basename=LOGFILEBASE     Base name for log files
  -V, --version                      Print version
\end{SaveVerbatim}
\respuesta{codigo}

A continuación hay que añadir el servicio swat al archivo \texttt{services} localizado en la carpeta \texttt{/etc}.

\begin{listing}[style=consola, numbers=none]
 # echo "swat	901/tcp" >> /etc/services
\end{listing}

Lo siguiente que hay que hacer es habilitar el uso de SWAT en xinetd. Para ello se necesita conocer el directorio donde se encuentra el binario. Normalmente será uno de los siguientes:
\begin{itemize}
 \item \texttt{/usr/local/samba/bin} la localización por defecto de Samba
 \item \texttt{/usr/sbin} localización por defecto en la mayor parte de los sistemas Linux
\item \texttt{/opt/samba/bin}
\end{itemize}

Se puede comprobar la localización del binario usando el comando \emph{whereis}.

\begin{listing}[style=consola, numbers=none]
$ whereis swat
\end{listing}

\begin{SaveVerbatim}{codigo}
swat: /usr/sbin/swat /usr/share/man/man8/swat.8.gz
\end{SaveVerbatim}
\respuesta{codigo}

Una vez localizado el directorio donde se encuentra el binario de SWAT (en este caso es \texttt{/usr/sbin}),hay que editar los archivos de configuración de xinetd, que suelen ser \texttt{xinetd.conf} y otros archivos localizados en la carpeta \texttt{xinetd.d} dentro del directorio \texttt{/etc}.

El archivo \texttt{xinetd.conf} es el archivo general de configuración. Su contenido por defecto es el siguiente:

\begin{listing}
  Simple configuration file for xinetd
#
# Some defaults, and include /etc/xinetd.d/

defaults
{

# Please note that you need a log_type line to be able to use log_on_success
# and log_on_failure. The default is the following :
# log_type = SYSLOG daemon info

}

includedir /etc/xinetd.d
\end{listing}

Al final del archivo se hace una inclusión de los archivos de la carpeta \texttt{/etc/xinetd.d}, que es donde se guardan como archivos individuales los datos de configuración de los diferentes servicios.

Para habilitar SWAT,  se crea un nuevo archivo llamado \texttt{swat} dentro del directorio \texttt{/etc/xinetd.d} con el siguiente contenido:

\begin{listing}
# default: off
# description: SWAT is the Samba Web Admin Tool. Use swat \
#              to configure your Samba server. To use SWAT, \
#              connect to port 901 with your favorite web browser.
service swat
{
	port    = 901
	protocol = tcp
	socket_type     = stream
	wait    = no
	only_from = localhost
	user    = root
	server  = /usr/sbin/swat
	log_on_failure  += USERID
	disable = no
}
\end{listing}

Nótese que en al campo ``server'' se le ha asignado la localización del binario de swat.


Para que el demonio xinetd relea su configuración hay que ejecutar lo siguiente:

\begin{listing}[style=consola, numbers=none]
# /usr/bin/killall --verbose -HUP xinetd
\end{listing}

\begin{SaveVerbatim}{codigo}
xinetd(7015) eliminado con la señal 1
\end{SaveVerbatim}
\respuesta{codigo}

Una vez hecho esto, el servicio SWAT queda activado y será accesible desde un navegador accediendo a localhost por el puerto 901:

\begin{listing}[numbers=none]
http://localhost:901
\end{listing}

Para entrar a la pantalla de configuración de SWAT hay que introducir el nombre de usuario y contraseña de root del sistema.

\begin{figure}[hbt!]
 \centering
 \includegraphics[width=0.95\textwidth]{Utils/swat_login.png}
 \caption{Pantalla de acceso a SWAT}
 \label{fig:swat_entrada}
\end{figure}

\newpage
\subsubsection{Resumen instalación}

\begin{enumerate}
 \item  Instalación de paquetes:
	\begin{listing}[style=consola, numbers=none]
# aptitude install samba xinetd swat
	\end{listing}

\item Añadir swat a la lista de servicios:
	\begin{listing}[style=consola, numbers=none]
# echo "swat		901/tcp" >> /etc/services
	\end{listing}

\item Crear nuevo archivo de servicio para swat
\begin{listing}[style=consola, numbers=none]
# nano /etc/xinetd.d/swat
\end{listing}

\begin{listing}[style=archivo]
# default: off
# description: SWAT is the Samba Web Admin Tool. Use swat \
#              to configure your Samba server. To use SWAT, \
#              connect to port 901 with your favorite web browser.
service swat
{
	port    = 901
	protocol = tcp
	socket_type     = stream
	wait    = no
	only_from = localhost
	user    = root
	server  = /usr/sbin/swat
	log_on_failure  += USERID
	disable = no
}
\end{listing}

\end{enumerate}

\clearpage
\newpage
\section{Adaptación de ficheros a DokuWiki}
\label{sanity}

\subsection{Introducción}

En la sección \ref{Samba}. \nameref{Samba} de este documento se explicó como compartir en red algunas carpetas del servidor. Concretamente se establecieron como carpetas compartidas las que guardan las páginas y los archivos subidos al Dokuwiki, \texttt{/data/pages} y \texttt{/data/media} respectivamente.

Sin embargo, para evitar problemas de acceso a las páginas y los archivos creados de forma manual, éstos deben pertenecer al mismo usuario que si hubieran sido creados a través del DokuWiki, es decir, \textit{www-data}.

Adicionalmente, para que los archivos contenidos en la carpeta \texttt{/media} sean enlazables desde el DokuWiki, sus nombres de archivo no deben contener caracteres no soportados, como espacios en blanco o letras mayúsculas.

Para cumplir con estos requerimientos se automatizará la ejecución de un script que cambiará los nombres de archivo por otros más convenientes y cambiará el usuario al que pertenecen, siempre que esto sea necesario.

\subsection{Script para modificar nombres de ficheros}

Para modificar los nombres de los ficheros se utilizará el script \textit{sanity.pl}, escrito en Perl por Andreas Gohr, desarrollador principal de DokuWiki, y distribuido bajo una licencia libre.

El script es recursivo, y modifica los nombres de los archivos situados en la carpeta indicada convirtiendo todos los caracteres a minúsculas, sustituyendo los espacios en blanco por barras bajas, eliminando caracteres no permitidos y evitando que los guiones estén rodeados de rayas bajas. De esta forma DokuWiki no tendrá problemas para leer los ficheros subidos.

Se puede consultar el código del script en la sección \ref{script-sanity}. \nameref{script-sanity}.

\subsection{Script de shell}

Para posteriormente poder automatizar el proceso de ejecución del script de Perl \textit{sanity.pl} y la modificación del dueño de los archivos se creará el script de shell \texttt{sanity.sh}.

\begin{listing}[style=consola, numbers=none]
 # nano /usr/bin/sanity.sh
\end{listing}

\begin{listing}
#!/bin/bash
cd /var/www/dokuwiki/data 
chown -R www-data media
chown -R www-data pages
perl sanity.pl -l media
\end{listing}

Este script considera que \textit{sanity.pl} se encuentra en la carpeta en que se guardan los datos del DokuWiki, es decir \texttt{/var/www/dokuwiki/data}.

Antes de poder ejecutar el script habrá que hacerlo ejecutable.

\begin{listing}[style=consola, numbers=none]
 # chmod +x /usr/bin/sanity.sh
\end{listing}

Una vez hecho esto ya se podrá ejecutar el script escribiendo lo siguiente en la consola:

\begin{listing}[style=consola, numbers=none]
# sanity.sh
\end{listing}

El resultado de la ejecución de este script cuando hay archivos nuevos pendientes de modificar el nombre sería similar a lo siguiente:

\begin{SaveVerbatim}{codigo}
Renaming '5- Manual de usuario.pdf' to '5-manual_de_usuario.pdf'        OKAY
Renaming '1- Memoria.pdf' to '1-memoria.pdf'    OKAY
Renaming '2- Presupuesto.pdf' to '2-presupuesto.pdf'    OKAY
Renaming 'Jorge Cuervo Ortega' to 'jorge_cuervo_ortega' OKAY
Renaming 'FA_213_dDa.txt' to 'fa_213_dda.txt'   OKAY
Renaming 'FV_171_dDa.txt' to 'fv_171_dda.txt'   OKAY
Renaming 'FV_236_dDa.txt' to 'fv_236_dda.txt'   OKAY
Renaming 'FV_233_dDa.txt' to 'fv_233_dda.txt'   OKAY
Renaming 'filtrado_ECG.m' to 'filtrado_ecg.m'   OKAY
Renaming 'FV_241_2_dDa.txt' to 'fv_241_2_dda.txt'       OKAY
Renaming 'FV_246_dDa.txt' to 'fv_246_dda.txt'   OKAY
Renaming 'FV_234_dDa.txt' to 'fv_234_dda.txt'   OKAY
\end{SaveVerbatim}
\respuesta{codigo}

\subsection{Automatización del proceso}

Para que el script descrito en la sección anterior se ejecute de forma automática habrá que crear un archivo en \texttt{/etc/cron.d}, llamado por ejemplo \texttt{sanity}.

A continuación se muestra un ejemplo del contenido de este fichero para que el script se ejecute cada dos minutos de ocho de la mañana a diez de la noche y de lunes a viernes.

\begin{listing}[style=consola, numbers=none]
# nano /etc/cron.d/sanity
\end{listing}

\begin{listing}
#min hora dia mes dia_semana usuario comando
*/2 08-22 * * 1-5 root /usr/bin/sanity.sh
\end{listing}

\newpage
\subsection{Resumen}

\begin{enumerate}
 \item  Creación del script \texttt{sanity.pl} (ver código en \ref{script-sanity}. \nameref{script-sanity})
\begin{listing}[style=consola, numbers=none]
# nano /var/www/dokuwiki/data/sanity.pl
\end{listing}

\begin{listing}[style=consola, numbers=none]
 # chmod a+x /usr/bin/sanity.sh
\end{listing}

\item Automatización de la ejecución del script

\begin{listing}[style=consola, numbers=none]
 # nano /etc/cron.d/sanity
\end{listing}

\begin{listing}[style=archivo, numbers=none]
 #min hora dia mes dia_semana usuario comando
*/10 08-22 * * 1-5 root /usr/bin/sanity.sh
\end{listing}

\end{enumerate}
\newpage
\section{Copias de Seguridad}
\label{backups}

\subsection{Introducción}

Para prevenir posibles pérdidas de datos conviene disponer de un sistema de copias de seguridad. En las siguientes secciones se explica cómo se puede utilizar un disco duro compartido en red en el laboratorio del GSC para realizar de forma automática respaldos de los datos del DokuWiki.

\subsection{Automatización de las copias de seguridad}

Cron es un demonio que sirve para ejecutar tareas programadas según una combinación de la hora, día del mes, día de la semana y semana. Para su funcionamiento, Cron asume que el sistema estará encendido de forma continua. Si el sistema no está activo cuando está programada la ejecución de una tarea, ésta no se ejecutará \cite{cron}.

Si se quiere crear una tarea que se ejecute según una programación distinta a la hora, día, semana o mes, habrá que agregar un archivo que describa la programación en el directorio \texttt{/etc/cron.d}

\begin{listing}[style=consola, numbers=none]
# nano /etc/cron.d/backup
\end{listing}

El fichero debe tener el siguiente formato:

minuto   hora   día   mes   día\_de\_la\_semana usuario  comando

\begin{itemize}
    \item minuto — número entero entre 0 y 59
    \item hora — número entero entre 0 y 23
    \item día — número entero entre 1 y 31 (debe ser un día válido si se especifica un mes)
    \item mes — número entero entre 1 y 12 (o nombre corto del mes, por ejemplo, ene, feb, etc.)
    \item día\_de\_la\_semana — número entero entre 0 y 7 (donde 0 o 7 corresponde a Domingo), o el nombre corto del día de la semana, por ejemplo, dom, lun, etc.
    \item usuario — usuario que ejecutará el script
    \item comando — el comando a ejecutar
\end{itemize}

Para que el respaldo básico se ejecute de lunes a viernes y el total los sábados, el fichero \texttt{/etc/cron.d/backup} debe tener la siguiente forma:

\begin{listing}
# Backups basicos dokuwiki
00  22 * * 1,2,3,4,5    root        /usr/bin/backup.sh
#El backup basico se ejecuta a las 22:00 de lunes a viernes
# Backup total del Dokuwiki
00 22  * * 6         root /usr/bin/totalbackup.sh
#El backup total se ejecuta a las 22:00 de los sabados
\end{listing}

\newpage

% CONFIG: Estilo de los ap ndices
%\renewcommand\chaptername{Ap\'{e}ndice}                  % a partir de aqu�, el nombre del cap�tulo N sera "Ap�ndice N"
%\renewcommand\thechapter{\Roman{chapter}}                % cambiar el n�mero del cap�tulo a romanos
\renewcommand\thesection{\Roman{section}} % que las secciones sean del tipo "I.a", en vez de "1.1"
\renewcommand{\thesubsection}{\thesection.\alph{subsection}}
\setcounter{section}{0}                                  % empezar a contar desde 1 los "cap�tulos"

% CONTENIDO: Ap ndices, si se desean
\newpage
% \input{ListadoSc7ripts.tex}

% CONFIG: Estilo de la bibliograf a
\cleardoublepage                                  % empezar en p�gina a la derecha
\addcontentsline{toc}{section}{REFERENCIAS} % as� aparecer� en el ToC principal como "Bibliograf�a"
\bibliographystyle{unsrt}                          % el estilo de bibliograf�a. unsrt: números en orden, plain: números, alpha para nombre del autor y año
% CONTENIDO: Bibliografia. Introduce una lista de ficheros .bib que quieras usar, como se expone.
%\section{REFERENCIAS}
\bibliography{file-bibliography,file-bibliography2}
% CONTENIDO: Lista de publicaciones:
%\input{file-publications.tex}
\end{document}